\begin{exercise}[Non-uniqueness of risk-neutral probabilities in a trinomial tree]
Show that in a trinomial one-period model the risk-neutral probability need
not be unique.
\end{exercise}

\begin{solution}
In the binomial model, the two unknown probabilities are pinned down by two
linear equations, which is why \(\polyref{Thm.~1.2}\) gives a unique
risk-neutral probability.  In a trinomial one-step model, suppose the risky
return is one of
\[
  d<m<u
\]
and write
\[
  p_d=\Q(Y=d),\qquad p_m=\Q(Y=m),\qquad p_u=\Q(Y=u).
\]
The risk-neutral martingale condition for the discounted stock is
\[
  \EQ[Y]=e^{r\Delta t}.
\]
Thus the probabilities must solve
\[
  p_d+p_m+p_u=1,\qquad
  dp_d+mp_m+up_u=e^{r\Delta t}.
\]
This is a system of two equations in three unknowns.  If
\(a=e^{r\Delta t}\), then for any chosen value \(\lambda=p_m\),
\[
  p_u=\frac{a-d-\lambda(m-d)}{u-d},
  \qquad
  p_d=\frac{u-a-\lambda(u-m)}{u-d}.
\]
For an interval of values of \(\lambda\), these three numbers are non-negative
and sum to \(1\).  Hence there is generally a whole family of risk-neutral
probabilities.

For example, take
\[
  d=0.8,\qquad m=1,\qquad u=1.2,\qquad e^{r\Delta t}=1.05.
\]
Then
\[
  p_u=0.625-0.5\lambda,\qquad
  p_d=0.375-0.5\lambda,\qquad
  p_m=\lambda.
\]
Every \(\lambda\in[0,0.75]\) gives a valid risk-neutral probability.  The
risk-neutral probability is therefore not unique.  In financial terms, the
single risky asset and the money-market account do not span all three states
of the world in one step, so the one-step trinomial market is typically
incomplete; compare with \(\polyref{Def.~1.7}\).
\end{solution}
